%%%%%%%%%%%%%%%%%%%%%%%%%%%%%%%%%%%%%%%%%%%%%%%%%%%%%%%
% % Lines starting with % are comments, which are ignored.
% % This is a handy way of indicating the date and version of
% % your document, to wit:
% %
% % LaTeX sample file
% % Modified March, 2002
% %
%%%%%%%%%%%%%%%%%%%%%%%%%%%%%%%%%%%%%%%%%%%%%%%%%%%%%%%
\DocumentMetadata{
  % PDF standards: ONE key whose value is a comma-separated list --
  % UA-2 (accessibility) plus A-4f (archival; the archival companion of
  % UA-2). Do NOT use separate pdfstandard=... keys: with ordinary
  % key-value parsing only the last one takes effect.
  pdfstandard={UA-2,A-4f},
  lang=en-US,
  tagging=on,
  tagging-setup={
    math/setup=mathml-SE
  }
}

\documentclass[12pt]{article}
% homework-macros.tex -- shared package setup for homework-template.tex
% (article class).
%
% NOTE ON PATHS: every template in this repository is stored one level up
% from where it is compiled (see README.md). The \input path below is
% resolved against the folder the document is compiled in, which sits
% one directory above this file.

% amsmath + amssymb + amsthm: full math support and proof environments
% (\newtheorem, \begin{proof}, ...).
\usepackage{amsmath, amssymb, amsthm}

% unicode-math: math typeset with OpenType math fonts (needed under
% XeLaTeX/LuaLaTeX so math matches the text fonts).
\usepackage{unicode-math}

% hyperref: clickable links (the template uses \href for the textbook
% problem sources); loaded after the math packages so it can patch them.
% Uses hyperref's default link colors.
\usepackage[colorlinks=true]{hyperref}

% Shared course math shorthands (\RR, \NN, \ZZ, \QQ, \grad, \curl, \tr,
% \der, \pder, ...) from the repository root.
\newcommand{\RR}{\mathbb{R}}
\newcommand{\NN}{\mathbb{N}}
\newcommand{\ZZ}{\mathbb{Z}}
\newcommand{\QQ}{\mathbb{Q}}
\newcommand{\vv}{\vec{v}}
\newcommand{\ww}{\vec{w}}
\newcommand{\xx}{\vec{x}}
\newcommand{\yy}{\vec{y}}
\newcommand{\zz}{\vec{z}}
\newcommand{\uu}{\vec{u}}
\newcommand{\basis}[1]{\vec{e_{#1}}}
\newcommand{\id}{\mathrm{id}}
\newcommand{\CC}{\mathbb{C}}
\renewcommand{\div}{\operatorname{div}}
\newcommand{\grad}{\operatorname{grad}}
\newcommand{\curl}{\operatorname{curl}}
\newcommand{\tr}{\operatorname{tr}}
\newcommand{\der}[2]{\dfrac{d #1}{d #2}}
\newcommand{\pder}[3][]{\dfrac{\partial^{#1} #2}{\partial #3}}


% % Title and author(s)
%%%%%%%%%%%%%%%%%%%%%%%%%%%%%%%%%%%%%%%%%%%%%%%%%%%%%%%
\title{Homework 6: 09/24 Content}
\author{MATH 1271, Garrett Credi}
\date{09/24/2026}
%%%%%%%%%%%%%%%%%%%%%%%%%%%%%%%%%%%%%%%%%%%%%%%%%%%%%%%
% %
% % The next command allows your in import encapsulated
% % postscript files, .epsf or .eps files, which
% % contain vector graphic image data.
% %
%%%%%%%%%%%%%%%%%%%%%%%%%%%%%%%%%%%%%%%%%%%%%%%%%%%%%%%
% %   The first thanks indicates your affiliation
% %
% %  Just the name here.
% %
% % Your mailing address goes at the end.
% %
% % \thanks is also how you indicate grant support
% %
%%%%%%%%%%%%%%%%%%%%%%%%%%%%%%%%%%%%%%%%%%%%%%%%%%%%%%%

%%%%%%%%%%%%%%%%%%%%%%%%%%%%%%%%%%%%%%%%%%%%%%%%%%%%%%%
% $ The next command ensures the margins are 1 inch on
% % all sides.
% %
%%%%%%%%%%%%%%%%%%%%%%%%%%%%%%%%%%%%%%%%%%%%%%%%%%%%%%%
\begin{document}
\maketitle

Problems listed as AC are found in the \href{https://activecalculus.org/single1e/frontmatter.html}{Active Calculus} 
textbook at the end of the section listed.

Problems listed as APEX are found in the \href{https://opentext.uleth.ca/apex-calculus/apex-calculus.html}{APEX Calculus Textbook}
again at the end of the listed section.

Problems listed as PAULS-PRACT are found in \href{https://tutorial.math.lamar.edu/Problems/CalcI/CalcI.aspx}{Paul's Online Math Notes - Review Problems}
while problems listed as PAULS-ASSIGN are found in \href{https://tutorial.math.lamar.edu/ProblemsNS/CalcI/CalcI.aspx}{Paul's Online Math Notes - Assignment Problems}
and are indexed via [Topic]-[Section]-[Problem Number]

\section{Type 1 - Computing Derivative Functions}
Choose 4 of the following problems,
ones that you \textit{did not} choose
for the previous homework, to solve.

\begin{enumerate}
  \item PAULS-ASSIGN.Derivatives.The Definition of the Derivative.1-24 (each of these count as one problem)
\end{enumerate}

\section{Type 2 - Graphically Interpreting Derivative Functions}
Choose 3 of the following problems to solve.

\begin{enumerate}
  \item AC 1.4.1
  \item AC 1.4.2
  \item AC 1.4.5
  \item AC 1.4.10
  \item AC 1.4.12
\end{enumerate}

\section{Type 3 - Determining Information from Derivative Functions}
Choose 3 of the following problems to solve.

\begin{enumerate}
  \item PAULS-ASSIGN.Derivatives.Interpretation of the Derivative.13,20,21 (each of these counts as one problem)
  \item APEX.2.1.35
  \item APEX 2.1.36
\end{enumerate}

\section{Challenge Problem}
Compute the derivative functions for the following functions:
\begin{enumerate}
  \item $f(x) = x^1$
  \item $f(x) = x^2$
  \item $f(x) = x^3$
  \item $f(x) = x^4$
\end{enumerate}
Do you notice a pattern with the derivative functions you find?

Use the above calculations to guess what the derivative function of
\textit{the power function}, $f(x) = x^n$ where $n > 0$ is a positive integer, is.
Then, using the limit definition of the derivative, try and show that your guess is actually the correct formula for $f'(x)$.
You will, in some way shape or form, need to expand the 
binomial power $(x + h)^n$ - think about how many terms you \textit{actually need} and how you might go about determining the coefficients of those necessary terms.

% Most people don't use these, so they are "commented out"
% by starting the lines with a "%"
%\begin{keywords}
%   \LaTeX, typesetting
%\end{keywords}

%\begin{AMS}
%   50C60, 18C25
%\end{AMS}

%%%%%%%%%%%%%%%%%%%%%%
% % Here is the start of the Text
%%%%%%%%%%%%%%%%%%%%%%

\end{document}
